% Cover letter for the submission to the Journal of Industrial Ecology.
% Build: pdflatex cover_letter.tex
\documentclass[10pt,a4paper]{letter}
\usepackage[utf8]{inputenc}
\usepackage[T1]{fontenc}
\usepackage[margin=2.5cm]{geometry}
\usepackage{hyperref}
\hypersetup{hidelinks}
\usepackage{enumitem}
\setlist{nosep}
\setlength{\parskip}{0.5em}
% left-align the closing and signature
\longindentation=0pt

\signature{Jakob Dürrwächter\\Corresponding author, on behalf of all authors}
\address{Jakob Dürrwächter\\Potsdam Institute for Climate Impact Research\\Potsdam, Germany\\\href{mailto:jakob.duerrwaechter@pik-potsdam.de}{jakob.duerrwaechter@pik-potsdam.de}}

\begin{document}

\begin{letter}{Editorial Office\\Journal of Industrial Ecology}

\opening{Dear Editor,}

The production of plastics, steel, and cement causes roughly one-fifth of global greenhouse gas emissions.
In addition to low-emission production technologies, circular economy measures are needed to reduce virgin material production.

Assessing circular economy measures and their interaction with climate mitigation in basic material production requires prospective, global, economy-wide material flow and demand scenarios that cover these three materials consistently.

\textbf{Contribution:} We present REMIND-MFA, a global, top-down, economy-wide, stock-driven, prospective material flow analysis model for plastics, steel, and cement. To our knowledge, it is the first model with this specific combination of properties for all three materials, at country-level resolution or arbitrary aggregations thereof. Its further methodological contributions are:
\begin{itemize}
	\item a three-step stock extrapolation algorithm that derives regional demand scenarios directly from the gross domestic product and population trajectories of the Shared Socioeconomic Pathways, with minimal region-specific assumptions;
	\item a trade extrapolation method that accounts for shifting trade patterns;
	\item historical data through at least 2022;
	\item an open-source, modular, automated code base designed for coupling with other models, in particular the integrated assessment model REMIND.
\end{itemize}
Along with this manuscript, we publish comprehensive scenarios of prospective stocks and flows with country-level resolution and a time horizon up to 2100, making them available for use by the research community.
In addition to the model code, the input data processing toolchain is openly available.

\textbf{Fit to the journal:} The paper's subjects -- material flow analysis, in-use stock dynamics, and the circular economy -- are central topics of the Journal of Industrial Ecology. The manuscript builds on and complements models published in this journal, such as MISO2 (Wiedenhofer et al., 2024) and ODYM-RECC (Pauliuk et al., 2021). The model and data are intended for research on prospective material demand, circularity scenarios, and the coupling of material flow analysis with integrated assessment or life-cycle assessment models.

Thank you for considering our manuscript.

\closing{Sincerely,}

\end{letter}
\end{document}
